\subsection{Моменты и жидкость}
%%%%%%%%%%%%%%%%%%%%%%%%%%%%%%%%%%%%%%%%%%%%%%%%%%%%
%%%%%%%%%%%%%%%%%%%%%%%%%%%%%%%%%%%%%%%%%%%%%%%%%%%%
%%%%%%%%%%%%%%%%%%%%%%%%%%%%%%%%%%%%%%%%%%%%%%%%%%%%
\z[{\begin{wrapfigure}[3]{r}{0.15\textwidth}%
		\vspace{-20pt}%
		\begin{flushright}%
 			\includegraphics[width=0.15\textwidth]{5.3.1}%
		\end{flushright}%
	\end{wrapfigure}%
	}]%
{Два одинаковых по форме цилиндрических поплавка удерживаются с помощью нитей погруженными в воду на половину своей длины. Определите отношение масс поплавков.}
%
{$\frac{m_1}{m_2}=3$}

%%%%%%%%%%%%%%%%%%%%%%%%%%%%%%%%%%%%%%%%%%%%%%%%%%%%
\z[{\begin{wrapfigure}[5]{r}{0.15\textwidth}%
		\vspace{-5pt}%
		\begin{flushright}%
 			\includegraphics[width=0.15\textwidth]{5.3.2}%
		\end{flushright}%
	\end{wrapfigure}%
	}]%
	{Тонкая однородная деревянная рейка одним концом шарнирно прикреплена к стенке сосуда с водой, а другим погружена в воду на 1/5  своей длины. Найдите плотность дерева, из которого изготовлена рейка, если плотность воды $\rho_0$. Куда направлена и чему равна сила реакции, возникающая в шарнире, если масса рейки равна $m$? 
	}
%
{\puan $ \rho_x = \frac{9}{25}\rho_0;$~\puan $F=\frac{4}{9}mg, $ вверх}

%%%%%%%%%%%%%%%%%%%%%%%%%%%%%%%%%%%%%%%%%%%%%%%%%%%%
\z[{\begin{wrapfigure}[6]{r}{0.2\textwidth}%
		\vspace{-15pt}%
		\begin{flushright}%
 			\includegraphics[width=0.2\textwidth]{5.3.3}%
		\end{flushright}%
	\end{wrapfigure}%
	}]%
{На лёгком рычаге  уравновешены два цилиндра, имеющие одинаковые размеры. При этом точка опоры делит рычаг в отношении $1$ к $3$, а цилиндры погружены в жидкость  (левый — на треть, а правый — на две трети объёма). Плотность левого цилиндра $\rho_{1} = 4 \, \text{г/см}^3$, а правого -- $\rho_{2} = 2.2 \, \text{г/см}^3$. Определите плотность жидкости $\rho$.
}
%
{$\rho = \frac{3}{5}(3\rho_2-\rho_1)=1.65\, \text{г/см}^3 $}

%%%%%%%%%%%%%%%%%%%%%%%%%%%%%%%%%%%%%%%%%%%%%%%%%%%%
\z[{\begin{wrapfigure}[6]{r}{0.4\textwidth}%
		\vspace{-15pt}%
		\begin{flushright}%
 			\includegraphics[width=0.4\textwidth]{5.3.4}%
		\end{flushright}%
	\end{wrapfigure}%
	}]%
{Прямоугольный лёгкий сосуд с жидкостью массой $m$ помещён на рычаг массой $3m$. В жидкость опущено тело массой $2m$, не касающееся дна сосуда и удерживаемое нитью, перекинутой через блок. Какой массы $m_x$ груз необходимо подвесить к противоположному концу нити для равновесия всей системы?}
%
{$m_x=\frac{3}{2}m$}


%%%%%%%%%%%%%%%%%%%%%%%%%%%%%%%%%%%%%%%%%%%%%%%%%%%%
\z[{\begin{wrapfigure}[5]{r}{0.2\textwidth}%
		\vspace{-15pt}%
		\begin{flushright}%
 			\includegraphics[width=0.2\textwidth]{5.3.5}%
		\end{flushright}%
	\end{wrapfigure}%
	}]%
{Длинный однородный брусок с поперечным сечением в виде прямоугольника со сторонами $a \neq b$ подвешен на двух вертикальных нитях, прикреплённых к одному из рёбер, над сосудом,  в который наливают воду. Когда в сосуд налили некоторое количество воды, два ребра бруска оказались точно на поверхности воды (вид сбоку со стороны вышеупомянутого поперечного сечения  показан на рисунке). Найдите плотность материала, из которого сделан брусок (плотность воды $\rho$ известна).}
%
{$\rho_x=\frac{2}{3}\rho$}


%%%%%%%%%%%%%%%%%%%%%%%%%%%%%%%%%%%%%%%%%%%%%%%%%%%%
\z[{\begin{wrapfigure}[4]{r}{0.2\textwidth}%
		\vspace{-10pt}%
		\begin{flushright}%
 			\includegraphics[width=0.2\textwidth]{5.3.6}%
		\end{flushright}%
	\end{wrapfigure}%
	}]%
{Лёгкий цилиндрический сосуд с жидкостью стоит на двух симметричных опорах. Над одной из них внутри сосуда привязан к дну полностью погружённый в жидкость 
поплавок объёмом $V$ и плотностью $\rho$. Над другой опорой висит привязанный снаружи грузик такого же объёма $V$ и плотностью $3\rho$. Плотность жидкости в сосуде равна $\rho_0$ ($\rho < \rho_0 < 3\rho$). Найдите модуль разности сил реакции опор.}
%
{$N_1-N_2=(\rho_0-\rho)mg$}

%%%%%%%%%%%%%%%%%%%%%%%%%%%%%%%%%%%%%%%%%%%%%%%%%%%%
\z[{\begin{wrapfigure}[2]{r}{0.3\textwidth}%
		\vspace{-10pt}%
		\begin{flushright}%
 			\includegraphics[width=0.3\textwidth]{5.3.7}%
		\end{flushright}%
	\end{wrapfigure}%
	}]%
{На однородном рычаге уравновешена кювета с жидкостью и плавающим в ней бруском. Масса бруска равна $m$, масса кюветы вместе с жидкостью $3m$. Найдите массу рычага $M$.}
%
{$M=8m$}